% !TEX root = ../main.tex
% ---------------------------------------------------------------------
%  B7 -- PID-Regler und Anti-Reset-Windup (Skript Kap. 8/9; AS 7/8;
%        Folien Reglerdiskretisierung; Skript Digitale Regelung Kap. 4)
% ---------------------------------------------------------------------
\vlsec{cB7}{B7}{PID-Regler und Anti-Reset-Windup}
\label{sec:pid}

\SW{Regelkreis} Sollwert $r$, Regelgrösse $y$, Regeldifferenz
$e=r-y$, Steuergrösse $u$, Störgrösse; Regler $C(s)$, Strecke $G(s)$.
Digital: AD-Wandlung \ra $e(k)$ \ra Regelalgorithmus \ra DA-Wandlung, alle
$T$ einmal.
\begin{ztab}{@{}lX@{}}
Anteil & Wirkung \\ \midrule
P ($K_p$) & sofort proportional; bleibender Fehler (ausser integrierende
  Strecke); zu gross \ra instabil \\
I ($T_i$) & summiert $e$ \ra kein bleibender Fehler; träge,
  Überschwingen, Windup \\
D ($T_d$) & reagiert auf Änderungsrate, dämpft; verstärkt Rauschen \ra
  immer mit Tiefpass ($N$) \\
\end{ztab}
\begin{formel}[Realer PID (Labor-Form)]
\[C(s)=\frac{U(s)}{E(s)}=K_p\Big(1+\frac{1}{T_is}+\frac{T_ds}{\frac{T_d}{N}s+1}\Big)\]
$T_i$ Nachstellzeit, $T_d$ Vorhaltzeit, $T_d/N$ Filterzeitkonstante
($N\approx1..100$, Labor 20). Umrechnung: $K_i=K_p/T_i$, $K_d=K_pT_d$.\\
\T{D auf $y$ statt $e$}: Sollwertsprünge würden einen Stoss erzeugen
\ra $u_d$ aus $-y$ berechnen (Minuszeichen!).
\end{formel}
\SW{Indirekter Entwurf} (Labor): analog entwerfen, dann diskretisieren
\ra keine neuen Methoden, intuitiv; aber nur PID und Abtastung nicht
explizit berücksichtigt. \T{Direkt}: Strecke $H(z)$, digital entwerfen.

\subsection{Diskretisierung der drei Anteile}
\label{rez:pid}
\begin{ztab}{@{}lX@{}}
P & $u_p(k)=K_p\,e(k)$ \\
I Rechteck rückw. & $u_i(k)=u_i(k-1)+\frac{K_pT}{T_i}\,e(k)$ \\
I Rechteck vorw. & $u_i(k)=u_i(k-1)+\frac{K_pT}{T_i}\,e(k-1)$ \\
\T{I Trapez} & $u_i(k)=u_i(k-1)+\frac{K_pT}{2T_i}\big(e(k)+e(k-1)\big)$ \\
D ungefiltert & $u_d(k)=\frac{K_pT_d}{T}\big(e(k)-e(k-1)\big)$ \\
\T{D gefiltert} & $u_d(k)=\frac{T_d}{T_d+NT}\,u_d(k-1)+\frac{K_pT_dN}{T_d+NT}\big(e(k)-e(k-1)\big)$ \\
D auf $y$ & wie oben mit $-\big(y(k)-y(k-1)\big)$ statt $\Delta e$ \\
Summe & $u(k)=u_p(k)+u_i(k)+u_d(k)$ \\
\end{ztab}
\begin{formel}[Herleitung D gefiltert (Rückwärts-Rechteck)]
$\big(\tfrac{T_d}{N}s+1\big)U_d=K_pT_d\,s\,E$, mit $s\to\frac{1-z^{-1}}{T}$:
\[\frac{T_d}{NT}\big(u_d(k)-u_d(k-1)\big)+u_d(k)=\frac{K_pT_d}{T}\big(e(k)-e(k-1)\big)\]
Mal $NT$: $(T_d+NT)\,u_d(k)=T_d\,u_d(k-1)+K_pT_dN\,\Delta e$
\ra nach $u_d(k)$ auflösen \ra Tabelle.\\
Summe \ra Rekursion: $\sum_{i=0}^{k}=\sum_{i=0}^{k-1}+$ neuer Summand
\ra konstante Rechenzeit und Speicher.
\end{formel}

\subsection{Virtuelle Strecke PT1T2 (AS 7)}
$G(s)=\frac{K}{(1+T_1s)(1+T_2s)}$, $a_1=e^{-T/T_1}$, $a_2=e^{-T/T_2}$:
\[H(z)=\frac{b_0z+b_1}{z^2+c_1z+c_2},\quad c_1=-(a_1+a_2),\; c_2=a_1a_2\]
\[b_0=\tfrac{K}{T_1-T_2}\big[T_1(1-a_1)-T_2(1-a_2)\big]\]
\[b_1=\tfrac{K}{T_1-T_2}\big[T_2a_1(1-a_2)-T_1a_2(1-a_1)\big]\]
\[y(k)=-c_1y(k-1)-c_2y(k-2)+b_0u(k-1)+b_1u(k-2)\]
\SW{Test} $K=10$, $T_1=0.1$, $T_2=0.01$; Regler $K_p=1/K$, $T_i=T_1$,
$T_d=T_2$ (kürzt die Strecken-Pole).

\subsection{Init/Run-Muster}
\label{sec:initrun}
Konstante Koeffizienten nur \emph{einmal} im Init rechnen, im Run nur die
Differenzengleichung \ra minimale CPU-Last pro Zyklus.
\begin{lstlisting}
TYPE System_Modus_Enum : (System_Init, System_Run); END_TYPE
TYPE PID_Parameter_Strukt :
STRUCT
  Kp, Ti, Td : REAL;  N : UINT;   ARW : BOOL;
  umin, umax, Tr : REAL;
END_STRUCT
END_TYPE

FUNCTION_BLOCK PID_Regler
VAR_INPUT
  r, y : REAL;  Modus : System_Modus_Enum;
  Param : PID_Parameter_Strukt;  T : REAL;
END_VAR
VAR_OUTPUT u, usat : REAL; END_VAR
VAR
  e, e1, y1, up, ui, ud : REAL;
  ci, cd1, cd2 : REAL;               // vorberechnete Koeffizienten
END_VAR

CASE Modus OF
System_Init:
  ci  := Param.Kp * T / (2.0 * Param.Ti);
  cd1 := Param.Td / (Param.Td + UINT_TO_REAL(Param.N) * T);
  cd2 := Param.Kp * Param.Td * UINT_TO_REAL(Param.N)
         / (Param.Td + UINT_TO_REAL(Param.N) * T);
  ui := 0.0; ud := 0.0; e1 := 0.0; y1 := y;
System_Run:
  e  := r - y;
  up := Param.Kp * e;
  ud := cd1 * ud - cd2 * (y - y1);         // D auf y
  u  := up + ui + ud;                      // ui(k-1)!
  IF u > Param.umax THEN usat := Param.umax;    // Sättigung
  ELSIF u < Param.umin THEN usat := Param.umin;
  ELSE usat := u; END_IF
  ui := ui + ci * (e + e1);                // Trapez
  IF Param.ARW THEN                        // Back-Calculation
    ui := ui + T / Param.Tr * (usat - u);
  END_IF
  e1 := e;  y1 := y;
END_CASE
\end{lstlisting}

\subsection{Stellgrössenbegrenzung und Anti-Reset-Windup}
\label{sec:arw}
\SW{Sättigung} Jeder Aktor ist begrenzt (Ventil offen, Spannung):
$u_{\text{sat}}=u_{\max}$ wenn $u\ge u_{\max}$, $u_{\min}$ wenn
$u\le u_{\min}$, sonst $u$. Die Begrenzung gehört physikalisch zur
Strecke, der Regler muss $u_{\text{sat}}$ aber kennen.
\begin{formel}[Windup]
Ist $u$ lange am Anschlag, integriert der I-Anteil die Regeldifferenz
ungebremst weiter, obwohl mehr Stellgrösse nicht wirkt. Nach dem
Vorzeichenwechsel von $e$ muss er erst wieder abgebaut werden \ra starkes,
langes \T{Überschwingen}, träge Reaktion auf neuen Sollwert, eventuell
Zahlenüberlauf.
\end{formel}
\begin{ztab}{@{}lX@{}}
Massnahme & Prinzip \\ \midrule
Conditional Integration & I-Anteil anhalten, solange gesättigt (und $e$
  weiter in dieselbe Richtung zeigt). Einfach, dynamisch nicht optimal. \\
\T{Back-Calculation} & Differenz $u_{\text{sat}}-u$ mit $\frac{1}{T_r}$
  gewichtet vor den Integrator zurückführen. Ohne Sättigung ist die
  Differenz 0 \ra kein Einfluss. \\
\end{ztab}
\begin{formel}[ARW diskret (Folien)]
$u_{\text{nosat}}(k)=u_p(k)+u_i(k-1)+u_d(k)$ \ra begrenzen \ra
\[u_i(k)=u_i(k-1)+\frac{K_pT}{T_i}\,\frac{e(k)+e(k-1)}{2}
+\frac{T}{T_r}\big(u_{\text{sat}}(k)-u_{\text{nosat}}(k)\big)\]
Startwert $T_r=\sqrt{T_iT_d}$ oder $\frac{T_i+T_d}{2}$ (Labor
Motor: $T_r=\sqrt{T_iT_d}$, $N=20$).
\end{formel}
\SW{Test AS 8} $u_{\min}=0$, $u_{\max}=10$; Sollwert als PWM 25 \ra 75;
$K_p\in\{\frac1K,\frac{10}{K},\frac{100}{K}\}$ -- hohe Verstärkung zeigt
den Windup am deutlichsten; wird es instabil, $T$ verkleinern.

\subsection{Laborversuch (Semesterleistung)}
SPS: Messsignale im Scope, Aktoren im Config-Modus, Zustandsmaschine
Hand/Auto bzw.\ Open-/Closed-Loop in \R{MAIN}, PID-FB, Prozess-FB.
Identifikation per \T{Sprungantwort} ($K,\tau_1,\tau_2$), Entwurf
(Ziegler-Nichols, SISO-Tool), Reserven: Phase $\ge50^\circ$,
Amplitude $\ge8$\,dB. Encoder EL5101 zählt alle Flanken von A und B
\ra $4\times$ Auflösung (512 \ra 2048/Umdrehung); Zählerüberlauf abfangen.
Motor EL7342: $\pm32767\leftrightarrow\pm24$\,V.
